% =====================================================================
\section{L'analogia con ACTN (RFC 8453)}
\label{sec:actn}
% =====================================================================

La GS~QKD~018 non inventa un modello nuovo: applica alla QKD un pattern già
standard nelle reti di trasporto, il framework ACTN
(\emph{Abstraction and Control of TE Networks}) dell'IETF~\cite{rfc8453}.
ACTN prevede tre livelli: il \emph{Customer Network Controller} (CNC), il
\emph{Multi-Domain Service Coordinator} (MDSC) e i \emph{Provisioning
Network Controller} (PNC) di ciascun dominio, collegati dalle interfacce
CMI (CNC–MDSC) e MPI (MDSC–PNC). La corrispondenza con la GS~QKD~018 è
diretta (Tab.~\ref{tab:actn}).

\begin{table}[h]
\centering\small
\caption{Corrispondenza tra ACTN e la pila ETSI per la QKD.}
\label{tab:actn}
\begin{tabularx}{\textwidth}{@{}l l L@{}}
\toprule
\textbf{ACTN} & \textbf{ETSI QKD} & \textbf{Osservazione} \\
\midrule
CNC & amministratore / portale cliente & origina la richiesta di servizio \\
MDSC & SDN orchestrator & vista multi-dominio, decisione del cammino \\
PNC & SDN controller del dominio QKD & calcola i candidati, attiva i cammini \\
MPI & interfaccia GS~QKD~018 & topologia astratta, provisioning, notifiche \\
SBI & interfaccia GS~QKD~015 & controllo dei singoli nodi \\
TE topology astratta & \texttt{etsi-qkd-sdn-topology} & vista vendor-neutral \\
banda TE & SKR / ESKR & la risorsa da allocare è la chiave \\
\bottomrule
\end{tabularx}
\end{table}

La differenza sostanziale sta nella risorsa. In ACTN la banda di un link è
fissata dalla capacità del canale; in QKD l'SKR fluttua con le condizioni
fisiche e la chiave si accumula in buffer finiti, da cui viene consumata.
Per questo nella 018 notifiche come \texttt{\_perf\_update} e
\texttt{\_overloaded} e la politica di \texttt{path\_restoration} hanno un
ruolo più centrale che nelle reti di trasporto classiche.
